% Proved after four actual-file audits. Separate post-CP23 research; ready for integration.
\subsection{Quadratic counts on a curve with prescribed full local types}

\begin{lemma}[Proved: finite integral jets determine cyclic unit identifications]
\label{D:curve-integral-unit-jet-comparison}
Let $A$ be any algebra over a field of characteristic $p>0$.
Write $F$ for coefficientwise $p$-th power, including $t\mapsto t^p$,
on $A[[t]]$. Fix $f,r\geq1$ and cyclic families of matrices
$U_a,V_a\in\operatorname{Mat}_r(A[[t]])$, indexed by
$a\in\mathbb Z/f\mathbb Z$. Assume that they are invertible over
$A((t))$ and that $t^cU_a^{-1}$ and $t^cV_a^{-1}$ are integral,
for one fixed integer $c\geq0$.
Choose integers $N\geq1$ and $L=N+c$ with $(p-1)N>c$.
Suppose that invertible matrix jets $h_a$ modulo $t^L$ satisfy
\[
 h_{a+1}U_a-V_aF(h_a)=0\pmod {t^L}.
 \tag{D.CL.1}
\]
There is a unique family of integral invertible matrices $H_a$
congruent to $h_a$ modulo $t^N$ and satisfying
\[
                  H_{a+1}U_a=V_aF(H_a).
 \tag{D.CL.2}
\]
The construction commutes with every homomorphism of coefficient
algebras. Only the retained $N$-jet, not necessarily the whole
$L$-jet, is preserved. Conversely every integral unit identification
supplies jets satisfying (D.CL.1).
\end{lemma}

\begin{proof}
Lift the jets arbitrarily to integral invertible matrices $H^{(0)}_a$.
Such lifts exist since invertibility of a matrix over $A[[t]]$ is
equivalent to invertibility of its constant term. Define
\[
 (\mathcal T H)_{a+1}=V_aF(H_a)U_a^{-1}.
 \tag{D.CL.3}
\]
The first difference $\mathcal T H^{(0)}-H^{(0)}$ is divisible by
$t^{L-c}=t^N$. If two families differ by $t^j$, their images under
$\mathcal T$ differ by $t^{pj-c}$. Thus $\mathcal T$ preserves
the integral $t^N$-ball about $H^{(0)}$ and is contracting there,
since $pN-c>N$. Every iterate remains invertible by its unchanged
constant term. The $t$-adic limit exists in every coefficient ring
and satisfies (D.CL.2). The same estimate proves uniqueness in
that ball. Different arbitrary lifts of the prescribed jets have
the same limit, and uniqueness proves compatibility with base
change. No division by $p$, no flatness of Frobenius on $A$, and
no bound on the nilpotence of $A$ are used.
\end{proof}

\begin{theorem}[Proved: fixed full local linear types have quadratic growth]
\label{D:curve-fixed-local-linear-growth}
Let $X$ be a smooth proper connected curve over an algebraically
closed field $k$ of characteristic $p>0$, let $D$ be a nonempty
finite reduced divisor, and put $X^\circ=X\setminus D$.
Fix a real number $b\geq0$. There is a constant $C$, depending
only on $(X,D,b,p)$, with the following property.
For $q=p^f$, $r\geq1$, and each $x\in D$, fix a continuous
local representation
\[
 \tau_x:\operatorname{Gal}(\overline{k((t_x))}/k((t_x)))
                  \longrightarrow\operatorname{GL}_r(\mathbb F_q)
\]
up to conjugacy. Suppose that every quotient defined by $\tau_x$
has trivial upper ramification groups above $b$.
Then the number of conjugacy classes of continuous representations
$\pi_1(X^\circ)\to\operatorname{GL}_r(\mathbb F_q)$ whose full
local types are the prescribed $\tau_x$ is at most
\[
                              q^{Cr^2}.
 \tag{D.CL.4}
\]
The constant is independent of the local splitting extensions,
their degrees, the actual matrices of the local types, $r$ and $f$.
\end{theorem}

\begin{proof}
We give both the finite presentation and the integral local
comparison needed for the component count. Bounds for ramification
breaks alone are not used as a replacement for full local types.

\emph{Canonical lattices and bounded slopes.}
For a global representation in question, normalize $X$ in a finite
Galois cover trivializing it. Write its finite group as $G$ and its
coefficient twists as $V_a=V^{[p^a]}$, for $a\in\mathbb Z/f\mathbb Z$.
The coherent sheaves
\[
                 E_a=(\pi_*\mathcal O_{X'}\otimes_k V_a)^G
 \tag{D.CL.5}
\]
are torsion-free, hence vector bundles of rank $r$ on the regular
curve $X$. Here $V_a$ denotes the corresponding $k$-vector space
with its specified finite-field matrices. Invariants are kernels;
no exactness of $G$-invariants in characteristic $p$ is asserted.
At $x\in D$, their completions are exactly the invariant lattices
of Lemma~\ref{D:different-Frobenius-invariant-lattice} for $\tau_x$
and its coefficient twists. This follows by the decomposition of
the completed normalization into its local factors and taking
invariants in the induced module. In particular these local
lattices depend functorially only on the local representation.
Coordinatewise Frobenius gives maps
\[
 F_X^*E_a\xrightarrow{\phi_a}E_{a+1},\qquad
 E_{a+1}(-cD)\xrightarrow{\psi_a}F_X^*E_a,
 \quad c=\lceil p(b+1)\rceil,
 \tag{D.CL.6}
\]
whose two composites are multiplication by the canonical Cartier
section of $cD$. The different estimate in the cited lemma proves
these assertions at every marked completion; they are isomorphisms
on $X^\circ$, so they hold globally. The same conclusion includes
trivial local types.

For a vector bundle $E$ on a curve one has
$p\mu_{\max}(E)\leq\mu_{\max}(F_X^*E)$, by pulling back its
maximal-slope subbundle. An injection of bundles bounds the maximal
slope of the source by that of the target. Applying this first to
$\phi_a$ and then to the dual of $\psi_a$, and cycling the $f$
inequalities, gives
\[
 -\frac{c\deg D}{p-1}\leq\mu_{\min}(E_a)
                  \leq\mu_{\max}(E_a)\leq0.
 \tag{D.CL.7}
\]
This argument requires no semistability assertion for tensor
products in positive characteristic. Every fixed twist of these
bundles has $h^0+h^1=O(r)$, with constants depending only on the
fixed curve and the twist.

\emph{An actual polynomial atlas with fixed terms.}
We spell out the ordinary-curve specialization of the fixed
partial-resolution construction in
Lemma~\ref{D:fixed-partial-resolution-unit-atlas}. There is no
root stack, normal layer or moving determinant parameter here.
Choose a fixed line bundle $H$ of sufficiently large positive degree
so that $E_a^+=E_a\otimes H$ has minimum slope greater than
$2g(X)+2$ whenever (D.CL.7) holds. A general $(r+1)$-tuple of
sections generates $E_a^+$: the failure locus of a general such
tuple has expected codimension two on the curve. Its kernel is
a line $L_a^{-1}$, where $L_a=\det E_a^+$ is globally generated
and has $H^1=0$. Choose a fixed positive line bundle $B$ and
$b_r=\kappa r$ with one fixed integer $\kappa$ large enough that
$L_a^{-1}B^{b_r}$ is also globally generated with $H^1=0$ for every
allowed $E_a$. These choices are possible because $\deg L_a=O(r)$
and the lower bound on its degree is strictly above $2g(X)-2$.

Use the fixed terms
\[
 P_0=\mathcal O_X^{r+1},\qquad
 P_j=(B^{-j b_r})^{\oplus2}\quad(1\leq j\leq7).
 \tag{D.CL.8}
\]
Two generators at each later step give exact middle terms with
alternating line syzygies
$L_a^{-1}B^{-2v b_r}$ and $L_aB^{-2v b_r}$.
The generating line at an odd step is $L_a^{-1}B^{b_r}$ and
at an even step is $L_a$. Their $H^1$ vanishing also lifts any
specified presentation across a small extension of local Artinian
coefficient algebras. This proves coverage, including coverage of
such prescribed lifts, without choosing a basis of a moving Picard
variety.

Parameterize the actual global differentials between (D.CL.8),
impose $d^2=0$, and impose the open rank-one conditions on every
consecutive map. Over any coefficient algebra the unit-minor
cancellation proves that the middle cohomology is zero and the two
end cohomology modules are bundles. Thus the degree-zero cohomology
is an actual rank-$r$ bundle $E_a^+$ over the proper curve. The
rank conditions are open on the coefficient scheme because their
failure locus on the proper curve has closed image. The slope
condition (D.CL.7) is the proper-Quot open condition already proved
in the cited atlas. It introduces no equations in the ambient
coordinate count. Untwist by the fixed $H$ to obtain $E_a$.
All these assertions hold over arbitrary coefficient algebras;
rank is imposed by unit minors on the actual local rings.

Here are the exact Hom witnesses used for (D.CL.6). Truncate the
source presentation at length seven and the target at length three.
Their end-error triangles give error intervals $[-3,-2]$ and
$[4,8]$, disjoint from degrees zero, one and two. Consequently
finite derived Hom computes exactly the ordinary Hom of the
bundles in degree zero. Compose through the inclusion of the
length-three target in its length-seven presentation. To test an
identity, precompose with its surjective top term $P_0$; do not
impose an identity on an equal-length partial complex. A degree
minus-one boundary witness tests equality exactly. This is the
separated-length argument (D.FR.5)--(D.FR.7), applied on $X$.

For clarity, the finite complexes are obtained by splitting the
two-affine Cech complexes of the fixed line-bundle terms over $k$.
The only perturbations are the variable internal differentials.
Every perturbation word increases internal degree, whose width is
bounded by seven and by the fixed composition diagram. Hence the
homological perturbation formulas terminate after a fixed number
of factors. They compute cycles, boundaries and composition over
every coefficient algebra, by the explicit identities previously
proved in that construction and
\citep[Sections~2.1--2.4 and~3.2]{C_Crainic2004Perturbation}.
This is a finite algebraic construction, not a dimension estimate
for a deformation space.

Every top-to-top Hom block has multiplicity $O(r^2)$ and fixed line
degree. A block with one positive-index term has multiplicity
$O(r)$ and line degree $O(r)$; a block with two such terms has
bounded multiplicity and line degree $O(r)$. Their total $h^0+h^1$
is $O(r^2)$. Fixed twists, the poles $cD$, and one-step Frobenius
multiply these bounds only by a constant. Frobenius pulls back
locally split complexes exactly over every coefficient algebra,
without a flatness assumption on that algebra, and replaces
variable coefficients by their $p$-th powers. We therefore obtain
all $f$ objects, maps $\phi_a,\psi_a$, and their exact inverse
identities with $O(fr^2)$ variables and equations of degree at
most a constant $d_0$ depending only on the fixed data and $p$.

\emph{Full local types from fixed-length jets.}
Choose the same $N,L=N+c$ as in
Lemma~\ref{D:curve-integral-unit-jet-comparison}. For each $x$ and
phase fix a basis of the canonical local lattice of $\tau_x^{[p^a]}$,
and write its Frobenius matrix as $V_{x,a}(t_x)$. The actual
coefficients of these matrices are constants, however complicated
the splitting cover. Only their first $L$ coefficients occur in
the presentation. Their inverse-pole bound is the common $c$.

On $X\times\operatorname{Spec}A$ impose frames of $E_a$ on
$D_L=\coprod_x\operatorname{Spec}k[t_x]/(t_x^L)$, with inverse
witnesses, and impose (D.CL.1) for the frame matrices and
$V_{x,a}$. These add $O(fr^2)$ variables of degree bounded in
terms of $p$ and $L$. To make the frame construction explicit,
for a local presentation $P_1\xrightarrow d P_0\to E_a\to0$
on $D_L$, use matrices $u:P_0\to A[t]/(t^L)^r$ and
$v:A[t]/(t^L)^r\to P_0$ with
\[
 u d=0,\qquad uv=1,\qquad vu-1=d w.
 \tag{D.CL.9}
\]
The last equality is witnessed by a matrix $w:P_0\to P_1$.
The locally split middle exactness makes these precisely a frame
of the quotient with a chosen lift. All sizes are $O(r)$ and $L$
is fixed. Restriction of arrows uses the separated partial
complexes before taking their cohomology; their locally free end
terms remain so after restriction. The same finite boundary
witnesses therefore express the induced jet maps exactly. There
are no omitted negative Tor terms. Frames of Frobenius pullbacks
are obtained by coefficientwise $p$-th power; the induced maps
are polynomials of bounded degree.

Let $Z$ be the resulting open locus of the affine polynomial
witness scheme. Its underlying affine equations have at most
$A_0fr^2$ variables and degree at most $d_1$, independent of the
local types. Every geometric point gives actual bundles and maps
(there is no formal algebraization step on this proper curve).
The jet comparison lemma gives actual integral cyclic unit
identifications of their completed local lattices with the fixed
local models. On $X^\circ$, the $f$ cyclic arrows give a unit
$q$-Frobenius bundle. The fixed-frame equations give a finite etale
$\operatorname{GL}_r(\mathbb F_q)$-torsor, exactly as in the
previous unit-bundle comparisons. Conversely every representation
being counted has its canonical family (D.CL.5), lies in (D.CL.7),
and is represented in $Z$. The local contractions commute with
coefficient base change.

\emph{Why unit-isomorphism classes are open.}
Fix a represented unit object $E^0$. The functor of isomorphisms
from $E^0$ to the universal family on the proper curve, with all
inverse and cyclic unit equations, is of finite presentation.
Indeed proper Hom of flat coherent bundles is affine of finite
presentation, and the displayed finite equations cut out this
functor
\citep[Proposition~99.3.10, Tag~08K6, and the stated inverse equations]
{StacksUnitJetTransferD}.
We check the required lifting at a specified Artinian point, rather
than only its tangent space. Let $A'\to A$ be a nilpotent small
extension of finite local Artinian $k$-algebras and let a unit
isomorphism be specified over $A$. Nilpotent invariance of finite
etale covers lifts it uniquely on $X^\circ_{A'}$
\citep[Lemma~58.8.3, Tag~0BQB, complete proof]
{StacksUnitJetTransferD}.
At each puncture the retained jets identify both objects, integrally
and compatibly over $A'$ and $A$, with their fixed canonical local
unit lattices. After these identifications, the local reduction
of the specified isomorphism is an automorphism of the constant
local $\mathbb F_q$-system. Its unique constant lift preserves the
canonical invariant lattice: the invariant-lattice construction
(D.CL.5) is functorial for every intertwiner. Full faithfulness of
nilpotent etale invariance says that the restriction of the global
lift is exactly this integral local lift. It therefore has no
pole at any point of $D$; its inverse has none either. Faithful
flatness of completion, applied to the coherent Hom bundles,
extends the isomorphism uniquely over all of $X_{A'}$.

This proves lifting of every specified Artinian point for this
finite-presentation Isom functor. It is formally etale, hence
etale near these points, by the Artinian local criterion and the
finite-presentation etaleness criterion
\citep[Tags~02HT, 02HM and~00UR, complete written proofs]
{StacksUnitJetTransferD}. In particular its image contains an
open neighborhood of every point in the same unit class.
Any unit isomorphism on $X^\circ$ between represented objects
preserves their canonical completed lattices by the same local
functoriality, so it is the proper-curve isomorphism just used.
Thus each represented class has an open image in $Z$.
Distinct classes give disjoint open images: an isomorphism after
an algebraically closed coefficient extension would give a point
of a finite-type Isom scheme, hence already a $k$-point.
Two disjoint nonempty opens cannot meet the same irreducible
component. The number of classes is therefore at most the number
of irreducible components of $Z$.

The ordinary-weight specialization of
Lemma~\ref{D:weighted-affine-component-bound} bounds this number
by $(A_0fr^2+1)d_1^{A_0fr^2}$, since passing to an open locus
cannot increase the number of irreducible components. Enlarging
the constant gives (D.CL.4). The local identifications are essential
in this argument: bounded poles by themselves would not justify
a nilpotence-independent extension of a unit isomorphism.
\end{proof}

\begin{corollary}[Proved: fixed full local projective types have quadratic growth]
\label{D:curve-fixed-local-projective-growth}
Under the hypotheses on $(X,D,b)$ above, fix full local conjugacy
classes $\sigma_x$ with values in $\operatorname{PGL}_r(\mathbb F_q)$,
all with highest upper break at most $b$. The number of global
projective conjugacy classes with these types is at most $q^{C'r^2}$,
with $C'$ depending only on $(X,D,b,p)$. More generally, allowing
at most $q^{a_xr^2}$ types at the point $x$ gives at most
\[
                            q^{(C'+\sum_x a_x)r^2}
 \tag{D.CL.10}
\]
global projective classes.
\end{corollary}

\begin{proof}
Put $n=q-1$. Since $D$ is nonempty, $X^\circ$ is a smooth affine
curve. Its etale $H^2(X^\circ,\mu_n)$ vanishes
\citep[Lemma~59.69.3, Tag~03RR, complete proof]{StacksCurveFixedLocalD}.
The already established degree-two injection from fundamental-group
cohomology to etale cohomology therefore gives
$H^2(\pi_1(X^\circ),\mathbb F_q^*)=0$; see also the precise
Cartan--Leray comparison previously used with
\citep[Theorem~14.9 and Remark~14.10]{MilneLEC2013H2D}.
The obstruction to lifting a continuous projective representation
through $\operatorname{GL}_r(\mathbb F_q)$ is consequently zero.
Thus every such representation has a continuous linear lift over
the same finite field. For $q=2$ the scalar kernel is trivial.

If the fibre being counted is empty there is nothing to prove.
Otherwise one global projective representation and its linear lift
just constructed supply a local linear lift of each $\sigma_x$.
After using a lifted projective conjugating matrix, any other
linear lift differs from it by a continuous scalar character into
$\mathbb F_q^*$. Kummer theory gives exactly $n$ such characters:
$k((t))^*/k((t))^{*n}\simeq\mathbb Z/n$, since the residue field
is algebraically closed and the units are $n$-divisible. All are
tame. Further, the prime-to-$p$ central scalar kernel adds no wild
upper group: for positive upper indices the kernel lies in a
$p$-group, and upper numbering is compatible with quotients.
Every linear lift therefore still has highest break at most $b$.
There are at most $n^{\deg D}$ possible tuples of linear local
types. Apply the theorem to each tuple and map linear conjugacy
classes to projective conjugacy classes. The factor $n^{\deg D}$
is absorbed in $q^{C'r^2}$. Summing over the stated allowed local
types proves (D.CL.10). No injectivity on global Brauer groups,
fixed neutralizing cover, or bound on the scalar order is assumed.
\end{proof}
